\documentclass[../../../include/open-logic-section]{subfiles}

\begin{document}

\olfileid{sth}{cardinals}{classing}
\olsection{সসীম, \usetoken{S}{enumerable}, \usetoken{S}{nonenumerable}}

অঙ্কবাচক সংখ্যার পরিচয় পাওয়ার পরে তাদের বিভিন্ন ধরন
নিয়ে একটু আলোচনা করা দরকার। বিশেষ করে সসীম,
!!{enumerable} ও !!{nonenumerable} অঙ্কবাচক সংখ্যা।

প্রথম দুটি ফলে দেখা যাবে, সসীম অঙ্কবাচক সংখ্যাগুলি
ঠিক সসীম ক্রমসংখ্যাই। \olref[z][infinity-again]{defnomega}-এ
এগুলিকেই আমাদের \emph{স্বাভাবিক সংখ্যা} হিসেবে সংজ্ঞায়িত করেছিলাম:

\begin{prop}\ollabel{finitecardisoequal}
$n,m\in\omega$ হোক। তবে $n=m$ তখনই, যখন $\cardeq{n}{m}$।
\end{prop}

\begin{proof}
\emph{বাম থেকে ডান} দিকটি সরল। \emph{ডান থেকে বাম}
প্রমাণ করতে ধরি, $n\neq m$ হলেও $\cardeq{n}{m}$।
ত্রিবিভাজনে $n\in m$ অথবা $m\in n$; সাধারণত্ব না হারিয়ে
$n\in m$ ধরি। তাহলে $n\subsetneq m$ এবং
একটি !!a{bijection} $f\colon m\to n$ আছে। এতে $m$
দেদেকিন্ড-অসীম হয়, যা
\olref[z][infinity-again]{naturalnumbersarentinfinite}-এর
বিরোধী।
\end{proof}

\begin{cor}\ollabel{naturalsarecardinals}
$n\in\omega$ হলে $n$ একটি অঙ্কবাচক সংখ্যা।
\end{cor}

\begin{proof}
তাৎক্ষণিক।
\end{proof}
\noindent
এ থেকে অঙ্কবাচক সংখ্যাকে ``সসীম'' বা ``অসীম''
বলার কয়েকটি স্বাভাবিক অর্থ যে একই, তাও পাওয়া যায়:
\begin{thm}\ollabel{generalinfinitycharacter}যেকোনো সেট $A$-র জন্য নিচের দাবিগুলি সমতুল্য:
\begin{enumerate}
	% BN-SRC-453: the cardinality, not A itself, is outside omega.
	\item\ollabel{card:notinomega} $\card{A}\notin\omega$,
	অর্থাৎ $A$-র অঙ্কবাচকতা স্বাভাবিক সংখ্যা নয়;
	\item\ollabel{card:omegaplus} $\omega\leq\card{A}$;
	\item\ollabel{card:infinite} $A$ দেদেকিন্ড-অসীম।
\end{enumerate}
\end{thm}

\begin{proof}
\olref[ord-arithmetic][using-addition]{ordinfinitycharacter},
\olref[cardinals][cardsasords]{lem:CardinalsBehaveRight}
এবং \olref{naturalsarecardinals} থেকে।
\end{proof}

এর ভিত্তিতে \olref[sfr][][]{part}-এ অপেক্ষাকৃত
অনানুষ্ঠানিকভাবে ব্যবহার করা কিছু ধারণা এবার
\emph{সংজ্ঞায়িত} করতে পারি:

\begin{defn}\ollabel{defnfinite}
$\card{A}$ স্বাভাবিক সংখ্যা, অর্থাৎ $\card{A}\in\omega$,
হলেই $A$-কে \emph{সসীম} বলব। অন্যথায় $A$-কে
\emph{অসীম} বলব।
\end{defn}
\noindent
তবে মনে রেখো, সংজ্ঞাটি $\ZFC$-র প্রেক্ষাপটে দেওয়া।
প্রত্যেক সেটের অঙ্কবাচকতা আছে, এটি নিশ্চিত করতে
আমাদের সুক্রমবিন্যাস লেগেছে। সুক্রমবিন্যাস ছাড়া এমন
সেট থাকতে পারে, যা সসীমও নয়, দেদেকিন্ড-অসীমও নয়।
\olref[choice][]{chap}-এ বিষয়টি আবার দেখব। আপাতত
সুক্রমবিন্যাস ব্যবহার করেই চলি।

এবার সসীম থেকে অসীম অঙ্কবাচক সংখ্যায় আসি।
প্রথমে দুটি সহজ কথা:

\begin{cor}\ollabel{omegaisacardinal}
$\omega$ ক্ষুদ্রতম অসীম অঙ্কবাচক সংখ্যা।
\end{cor}

\begin{proof}
$\omega$ একটি অঙ্কবাচক সংখ্যা। কারণ $\omega$
দেদেকিন্ড-অসীম; কোনো $n\in\omega$-র সঙ্গে
$\cardeq{\omega}{n}$ হলে $n$-ও দেদেকিন্ড-অসীম হতো,
যা \olref[z][infinity-again]{naturalnumbersarentinfinite}-এর
বিরোধী। এখন সংজ্ঞা অনুযায়ী $\omega$ ক্ষুদ্রতম
অসীম অঙ্কবাচক সংখ্যা।
\end{proof}

\begin{cor}
প্রত্যেক অসীম অঙ্কবাচক সংখ্যা একটি সীমা-ক্রমসংখ্যা।
\end{cor}

\begin{proof}
$\alpha$ একটি অসীম উত্তরসূরি ক্রমসংখ্যা হোক, অর্থাৎ
কোনো $\beta$-র জন্য $\alpha=\beta\ordplus 1$।
% BN-SRC-454: a finite predecessor would have a finite successor;
% finitecardisoequal does not establish the cited implication.
এই পূর্বসূরি সসীম হলে তার উত্তরসূরিও সসীম হতো;
অতএব $\beta$ অসীম। তাই
\olref[ord-arithmetic][using-addition]{ordinfinitycharacter}
অনুযায়ী $\cardeq{\beta}{\beta\ordplus 1}$। এখন
\olref[cardinals][cardsasords]{lem:CardinalsBehaveRight} থেকে
$\card{\beta}=\card{\beta\ordplus 1}=\card{\alpha}$;
সুতরাং $\alpha\neq\card{\alpha}$।
\end{proof}

\olref[sfr][siz][enm-alt]{defn:enumerable}-এই বলেছিলাম,
অসীম সেটের মধ্যে !!{enumerable} ও !!{nonenumerable}
আলাদা করা যায়। সেই সংজ্ঞা থেকে স্বাভাবিকভাবেই পাই:

\begin{prop}
$A$ !!{enumerable} তখনই, যখন $\card{A}\leq\omega$;
আর $A$ !!{nonenumerable} তখনই, যখন $\omega<\card{A}$।
\end{prop}

\begin{proof}
ত্রিবিভাজনে দুটি দাবি সমতুল্য, তাই $A$ !!{enumerable}
তখনই, যখন $\card{A}\leq\omega$, শুধু এটিই দেখালেই চলে।
\emph{ডান থেকে বাম}: $\card{A}\leq\omega$ হলে
\olref[cardinals][cardsasords]{lem:CardinalsBehaveRight}
ও \olref{omegaisacardinal} থেকে $\cardle{A}{\omega}$।
\emph{বাম থেকে ডান}: $A$ !!{enumerable} ধরি।
\olref[sfr][siz][enm-alt]{defn:enumerable} অনুযায়ী
তিনটি ক্ষেত্র সম্ভব:
\begin{enumerate}
	\item $A=\emptyset$ হলে
	\olref{naturalsarecardinals} ও
	\olref[cardinals][cardsasords]{lem:CardinalsBehaveRight}
	থেকে $\card{A}=0\in\omega$।
	\item $\cardeq{n}{A}$ হলে
	\olref{naturalsarecardinals} এবং
	\olref[cardinals][cardsasords]{lem:CardinalsBehaveRight} থেকে
	$\card{A}=n\in\omega$।
	\item $\cardeq{\omega}{A}$ হলে
	\olref{omegaisacardinal} থেকে $\card{A}=\omega$।
\end{enumerate}
অতএব সব ক্ষেত্রেই $\card{A}\leq\omega$।
\end{proof}
\noindent
এখানে $\omega$-র বিশেষ স্থান আছে। গণনীয় ক্রমসংখ্যা
অনেক থাকলেও:

\begin{cor}
$\omega$-ই একমাত্র !!{enumerable} অসীম অঙ্কবাচক সংখ্যা।
\end{cor}

\begin{proof}
$\cardfont{a}$ একটি !!a{enumerable} অসীম অঙ্কবাচক
সংখ্যা হোক। $\cardfont{a}$ অসীম বলে $\omega\leq\cardfont{a}$।
আবার $\cardfont{a}$ !!a{enumerable} অঙ্কবাচক সংখ্যা বলে
$\cardfont{a}=\card{\cardfont{a}}\leq\omega$।
তাই ত্রিবিভাজনে $\cardfont{a}=\omega$।
\end{proof}

অঙ্কবাচক সংখ্যা অবশ্যই অসীমসংখ্যক। তাই প্রশ্ন হতে পারে:
\emph{কতগুলি অঙ্কবাচক সংখ্যা আছে?} নিচের ফলগুলি
দেখায়, প্রশ্নটিই হয়তো নতুন করে ভাবা দরকার।

\begin{prop}\ollabel{unioncardinalscardinal}
$X$-এর প্রত্যেক সদস্য অঙ্কবাচক সংখ্যা হলে
$\bigcup X$-ও অঙ্কবাচক সংখ্যা।
\end{prop}

\begin{proof}
$\bigcup X$ যে ক্রমসংখ্যা, তা সহজেই যাচাই করা যায়।
$\alpha\in\bigcup X$ একটি ক্রমসংখ্যা হোক। তখন কোনো
অঙ্কবাচক সংখ্যা $\cardfont{b}$-র জন্য
$\alpha\in\cardfont{b}\in X$। $\cardfont{b}$ অঙ্কবাচক সংখ্যা
বলে $\cardless{\alpha}{\cardfont{b}}$। আবার
$\cardfont{b}\subseteq\bigcup X$ বলে
$\cardle{\cardfont{b}}{\bigcup X}$; সুতরাং
$\cardneq{\alpha}{\bigcup X}$। সাধারণভাবে,
$\bigcup X$ একটি অঙ্কবাচক সংখ্যা।
\end{proof}

\begin{thm}\ollabel{lem:NoLargestCardinal}
বৃহত্তম অঙ্কবাচক সংখ্যা নেই।
\end{thm}

\begin{proof}
যেকোনো অঙ্কবাচক সংখ্যা $\cardfont{a}$-র জন্য
কান্টরের উপপাদ্য (\olref[sfr][siz][car]{thm:cantor})
ও \olref[cardinals][cardsasords]{lem:CardinalsExist} থেকে
$\cardfont{a}<\card{\Pow{\cardfont{a}}}$।
\end{proof}

\begin{thm}
সব অঙ্কবাচক সংখ্যার সেট নেই।
\end{thm}

\begin{proof}
বিরোধ প্রতিষ্ঠার জন্য ধরি,
$C=\Setabs{\cardfont{a}}{\cardfont{a}\text{ is a cardinal}}$।
\olref{unioncardinalscardinal} অনুযায়ী $\bigcup C$
একটি অঙ্কবাচক সংখ্যা। তাই
\olref{lem:NoLargestCardinal} থেকে
এই সংখ্যার চেয়ে বড় একটি অঙ্কবাচক সংখ্যা
$\cardfont{b}>\bigcup C$ আছে। সংজ্ঞা অনুযায়ী $\cardfont{b}\in C$;
ফলে $\cardfont{b}\subseteq\bigcup{C}$, অর্থাৎ
$\cardfont{b}\leq\bigcup C$। এটি বিরোধ।
\end{proof}

রাসেলের কূটাভাস এবং বুরালি--ফোর্তির কূটাভাস,
দুটির সঙ্গেই এই ফলটির তুলনা করো।

\end{document}
